\begin{afterword}
\summaryhead

The post opens with people who died after playing online games for days. It calls the products that do this superstimuli: things that match an evolved preference more strongly than anything in the ancestral environment. A candy bar matches taste buds tuned to sugar, salt and fat; a retouched billboard model, as in Dove's ``Evolution'' video, beats any real woman; a video game can be more engaging than life. Because firms compete, the post argues, each will make its product a little more tempting than the rival's, past the point of harm to the consumer.

Why not just say no? Human willpower, the author suggests, is a limited compromise struck by evolution; resisting temptation spends an exhaustible supply of mental energy, and may spend it faster against temptations no hunter-gatherer met. The post asks whether ads for superstimuli harm even those who resist, and notes both that government may not be able to fix the problem and that this does not mean nothing is wrong. It ends with Simon Funk's novel \textsc{After Life}, in which humanity is marketed out of existence, and suggests that falling birth rates may come from ever more tempting alternatives to children.

\responsehead

The post states a vivid idea. Markets can supply stimuli stronger than any our preferences were tuned for, and competition pushes them past the point of harm. Its opening cases are real: the three linked reports match three deaths after gaming sessions of fifty hours to a week. The idea itself comes from ethology. Tinbergen found birds preferring oversized dummy eggs to their own and named the effect; the post does not mention this origin.

Most of what follows is argued by illustration and hedged guesses. The 5\% and 50\% of the market paragraph are an example. The account of how willpower evolved is offered with ``seems'', ``perhaps'' and ``probably'', and no evidence. The one supermodel example has no source. These are the ordinary tools of an essay, and the post marks most of them as guesses.

The exception is its mechanism. ``Resisting any temptation takes conscious expenditure of an exhaustible supply of mental energy'' is stated flatly, as the reason we cannot ``just say no.'' The link is to an economic model that assumes this. The psychology behind it, ego depletion, later failed two large preregistered replications, in 2016 and 2021. The post was written before them. The failures do not show that resisting temptation costs nothing; they fail to find the depletion the post names as the cost. The next paragraph's suggestion, that superstimuli drain willpower faster, rests on the same model.

The collapse in the title is argued only in the last paragraph, with ``Perhaps'' and a novel, and the author says openly that this is an argument from fiction. Its closing quotation, ``the human species was simply marketed out of existence,'' does not appear in the text of the novel that its author offers for download; the nearest line is a character's question, ``So, we simply marketed the humans out of existence?''

In short: a vivid statement of an idea from ethology, applied to markets with real cases and marked guesses; the one mechanism stated as fact has since failed two large replications, and the collapse of the title rests on ``Perhaps'' and a novel.
\end{afterword}
